% Template version as of 6/27/2024
\documentclass[conference, 12pt]{IEEEtran} 
\IEEEoverridecommandlockouts
% The preceding line is only needed to identify funding in the first footnote. If that is unneeded, please comment it out.
\usepackage{cite}
\usepackage{amsmath,amssymb,amsfonts}
\usepackage{algpseudocode}
\usepackage{algorithm}
\usepackage{graphicx}
\usepackage{textcomp}
\usepackage{booktabs}
\usepackage{listings}
\usepackage{csquotes}
\usepackage{hyperref}
\usepackage{xcolor}
\usepackage{fontawesome}
\usepackage{cleveref}

\crefformat{figure}{Fig.~#2#1#3}

\begin{document}

\title{
Literature Review on 3D Body Shape Analysis \\
for Disease Prediction and Healthcare 
}

\author{
\IEEEauthorblockN{
Anonymous Submission
% Aizierjiang Aiersilan
}

% \IEEEauthorblockN{
% % 1\textsuperscript{st}
% Aizierjiang Aiersilan
% % \textsuperscript{\href{mailto:xxx@example.com}
% % {\faEnvelopeO}
% % }\
% }
% \IEEEauthorblockA{
% % \textit{Faculty of Science and Technology} \
% University of Example\
% City, State\
% example@xxx.edu
% }

% \and

% \IEEEauthorblockN{
% % 2\textsuperscript{nd}
% XXXXX XXXXX
% }
% \IEEEauthorblockA{
% % \textit{Faculty of xxxxxxx xxxxxx} \
% XXXXX XXXXX \
% City, Country\
% xxxxx@xxxxx.cn
% }

}
\maketitle

\begin{abstract}
The integration of 3D body scanning technologies with advanced machine learning and computer graphics methodologies has revolutionized non-invasive health assessment and disease prediction. This comprehensive literature review examines the current state-of-the-art in 3D body shape analysis for healthcare applications, synthesizing research from leading computer vision, machine learning, and biomedical engineering venues. We systematically analyze methodological approaches ranging from traditional geometric feature extraction to cutting-edge deep learning architectures, including convolutional neural networks, point cloud processing networks, and generative adversarial networks. The review encompasses critical applications such as body composition estimation, metabolic syndrome prediction, liver health assessment, and maternal-fetal health monitoring. Through rigorous examination of recent advances, we identify key technical challenges including data acquisition standardization, model interpretability, and cross-population generalizability. We conclude by outlining promising future directions in multimodal data fusion, real-time processing, and explainable AI for clinical deployment. This review serves as a comprehensive resource for researchers and practitioners working at the intersection of computer graphics, computer vision, and digital health.
\end{abstract}

\begin{IEEEkeywords}
3D body scanning, machine learning, computer graphics, health assessment, digital health, point cloud processing, deep learning
\end{IEEEkeywords}


%%%%%%%%% BODY TEXT %%%%%%%%%
\section{Introduction}
\label{sec:intro}
\setlength{\parindent}{0pt}

Accurate and accessible health assessment is fundamental for early disease detection, personalized medical interventions, and effective public health strategies. Traditional methods for body composition analysis and internal health diagnostics, such as Dual-Energy X-ray Absorptiometry (DXA), Magnetic Resonance Imaging (MRI), and Computed Tomography (CT), offer high precision but are often limited by significant costs, the use of ionizing radiation (for DXA and CT), and restricted accessibility~\cite{zheng2024d3bt}.







\section{Search Strategy}

Search Years:
2016-2026 (10 years)

Search Keywords:
"3D body shape" OR "3D anthropometry" OR "human body modeling" AND (healthcare OR "medical applications" OR "body composition" OR diagnostics) AND (analysis OR measurement OR modeling)






\section{Conclusion}

This research proposal outlines a transformative approach to autonomous driving by integrating ...

%%%%%%%%% REFERENCES %%%%%%%%%
\newpage
{\small
\bibliographystyle{IEEETran}
\bibliography{IEEEfull}
}
\end{document}